\chapter{Representation of Emphasis in Malayalam}
\label{chap:introduction}

Emphasis is a linguistic means of highlighting a particular word, phrase, or piece of information in an utterance so that it receives greater prominence or attention. Languages convey emphasis through different grammatical and syntactic constructions, as well as through prosodic features such as stress, pitch, and intonation. The choice of mechanism depends on the structure and grammatical properties of the language.

\section{Prosodic and Syntactic Representation of Emphasis}

In Malayalam, emphasis can be expressed through both prosodic and syntactic means. Asher and Kumari \citep{asher1997malayalam} note that Malayalam falls into the syllable-timed category and that, except in emphatic speech, stress is not associated with greater loudness or intensity. As Moag \citep{moag2002malayalam} observes, while English speakers commonly convey focus by phonologically stressing the highlighted constituent during speech, Malayalam predominantly relies on constituent reordering and the syntactic device of clefting. Cleft sentences are far more common and natural in Malayalam than in English, serving as the standard syntactic vehicle for both information focus and contrastive emphasis.


\section{Clefting}

Clefting is a syntactic construction used to give emphasis to a particular constituent of a sentence. In Malayalam, Asher and Kumari \citep{asher1997malayalam} describe clefting as a very common construction. The process of clefting involves modifying the structure of an ordinary sentence so that a selected constituent receives emphasis. In a verbal sentence, the original finite main verb is first converted into a nominalised form by adding the suffix \textit{-aṭŭ}. The copula \textit{āṇŭ} is then associated with the constituent that is to be emphasised, while the nominalised form of the original verb occurs in the remaining part of the construction. Thus, the constituent followed by \textit{āṇŭ} becomes the element receiving emphasis, while the rest of the sentence is retained in the cleft construction.The copula can be added to a major constituent of the sentence, such as an NP, AdvP, or PostP, and that
constituent then carries a certain degree of emphasis. 

\begin{enumerate}[label=\textbf{\arabic*.}]

    \item \textbf{Subject Focus:}\\
    \textbf{Neutral:} Raman gave Krishnan two books yesterday.\\
    \ml{രാമൻ ഇന്നലെ കൃഷ്ണന് രണ്ട് പുസ്തകം കൊടുത്തു.}\\
    \textit{rāman innale kṛṣṇan-ŭ raṇṭŭ pustakaṁ koṭuttu.}\\
    \textbf{Cleft:} \textbf{RAMAN} gave Krishnan two books yesterday.\\
    \ml{രാമനാണ് ഇന്നലെ കൃഷ്ണന് രണ്ട് പുസ്തകം കൊടുത്തത്.}\\
    \textit{rāman-āṇŭ innale kṛṣṇan-ŭ raṇṭŭ pustakaṁ koṭuttatŭ.}\\
    \textbf{Meaning:} ``It was Raman who gave Krishnan two books yesterday.''

    \item \textbf{Temporal Adverbial Focus:}\\
    \textbf{Neutral:} Raman gave Krishnan two books yesterday.\\
    \ml{രാമൻ ഇന്നലെ കൃഷ്ണന് രണ്ട് പുസ്തകം കൊടുത്തു.}\\
    \textit{rāman innale kṛṣṇan-ŭ raṇṭŭ pustakaṁ koṭuttu.}\\
    \textbf{Cleft:} \textbf{YESTERDAY} Raman gave Krishnan two books.\\
    \ml{ഇന്നലെയാണ് രാമൻ കൃഷ്ണന് രണ്ട് പുസ്തകം കൊടുത്തത്.}\\
    \textit{innale-y-āṇŭ rāman kṛṣṇan-ŭ raṇṭŭ pustakaṁ koṭuttatŭ.}\\
    \textbf{Meaning:} ``It was yesterday that Raman gave Krishnan two books.''

    \item \textbf{Indirect Object Focus:}\\
    \textbf{Neutral:} Raman gave Krishnan two books yesterday.\\
    \ml{രാമൻ ഇന്നലെ കൃഷ്ണന് രണ്ട് പുസ്തകം കൊടുത്തു.}\\
    \textit{rāman innale kṛṣṇan-ŭ raṇṭŭ pustakaṁ koṭuttu.}\\
    \textbf{Cleft:} Raman gave \textbf{KRISHNAN} two books yesterday.\\
    \ml{രാമൻ ഇന്നലെ കൃഷ്ണനാണ് രണ്ട് പുസ്തകം കൊടുത്തത്.}\\
    \textit{rāman innale kṛṣṇan-āṇŭ raṇṭŭ pustakaṁ koṭuttatŭ.}\\
    \textbf{Meaning:} ``It was to Krishnan that Raman gave two books yesterday.''

\end{enumerate}

These examples show that the same basic cleft construction can be used to place emphasis on different constituents of the sentence. In each verbal cleft, the focused constituent is associated with \textit{āṇŭ}, while the original finite verb is replaced by its nominalised form.

\section{Distinguishing Clefting from Other Uses of \textit{āṇŭ}}

Mohanan and Mohanan \citep{mohanan1999two} identify two functions of \textit{āṇŭ} in Malayalam. It functions as an ordinary equative copula when it occurs between two nominative nominal constituents, and it functions as a cleft marker in certain non-equative constructions. Therefore, the presence of \textit{āṇŭ} alone is not sufficient to identify a cleft construction. Its function must be determined from the syntactic structure in which it occurs.

Accordingly, a sentence containing \textit{āṇŭ} must be examined with respect to the position of \textit{āṇŭ}, the constituent associated with it, the case pattern of the constituents, the presence or absence of a nominalised verb, and the overall structure of the clause. These structural properties help distinguish the ordinary equative use of \textit{āṇŭ} from its use as a cleft marker and from its occurrence in other constructions, such as aspectual constructions.

\subsection{Equative constructions}
When both arguments of the copula are nominative, \textit{āṇŭ} functions as an ordinary equative copula and normally occurs at the end of the clause \citep{mohanan1999two}. For example:
\begin{quote}
\textit{aruṇ abhibhāṣakan-āṇŭ}\\
``Arun is a lawyer.''
\end{quote}
Here, \textit{āṇŭ} connects the two nominative nominal constituents and does not by itself indicate a cleft construction.

\subsection{Full cleft constructions}
In a full cleft, \textit{āṇŭ} occurs immediately after the constituent being focused, while the remaining elements of the sentence follow it. The original main verb is nominalised and occurs later in the clause \citep{asher1997malayalam,mohanan1999two}.
\begin{quote}
\textit{rāman-āṇŭ innale kṛṣṇan-ŭ raṇṭŭ pustakam koṭuttatŭ}\\
``It was Raman who gave Krishnan two books yesterday.''
\end{quote}
Here, \textit{rāman} is the focused constituent and \textit{āṇŭ} immediately follows it. The original verb \textit{koṭuttu} is represented by the nominalised form \textit{koṭuttatŭ}. The combination of the focused constituent, the position of \textit{āṇŭ}, and the nominalised verb identifies the construction as a full cleft.

\subsection{Movement}

In an ordinary equative construction, \textit{āṇŭ} occurs with the predicate at the end of the clause. For example:

\begin{quote}
\textit{aruṇ abhibhāṣakan-āṇŭ}\\
``Arun is a lawyer.''
\end{quote}

When emphasis is given to the predicate constituent, the constituent associated with \textit{āṇŭ} can occur in a different position, with \textit{āṇŭ} occurring directly with the constituent receiving focus. For example:

\begin{quote}
\textit{abhibhāṣakan-āṇŭ aruṇ}\\
``It is a lawyer that Arun is.''
\end{quote}

Thus, the relevant change is not simply the movement of the focused constituent by itself. Rather, the focused constituent occurs together with \textit{āṇŭ}, resulting in a different constituent order from the ordinary equative construction. Jayaseelan observes that, in Malayalam cleft constructions, the focused phrase and the copula can move together and occur in different positions within the sentence \citep{jayaseelan2023}.

\subsection{Reduced cleft constructions}
When a non-nominative constituent, such as a dative or locative constituent, is combined with \textit{āṇŭ} without an overt lexical verb, the construction can also receive a reduced cleft interpretation. Following Mohanan and Mohanan \citep{mohanan1999two}, \textit{āṇŭ} functions as a cleft marker in such constructions:
\begin{quote}
\textit{aanakkŭ pani-āṇŭ}\\
``What the elephant has is a fever.''
\end{quote}
Thus, the absence of a nominalised main verb does not by itself exclude a cleft analysis. The case pattern and the interpretation of the construction must also be considered.

\section{Cleft Construnction in Malayalam}
\label{sec:moag_focus_typology}

In his descriptive reference grammar, Moag \citep{moag2002malayalam} establishes a systematic framework for how emphasis and focus operate across equative sentences, verbal sentences, information questions, and verbal predicates. Syntactic focus in Malayalam is realized by associating the copular focus marker \ml{ആണ്} (\textit{āṇŭ}) immediately with the targeted constituent, accompanied by structural adjustments conditioned by the predicate type.

\subsection{Focus in an Equative Sentence (No Lexical Verb)}
In equative (copular) sentences containing no lexical verb, the normal, neutral constituent order places the subject initially and the predicate nominal or locative phrase with the copula \ml{ആണ്} at the clause-final position.

To shift emphasis or focus to another constituent—such as a locative phrase—the copula \ml{ആണ്}, which usually occurs at the end of the sentence, is moved to the left to immediately follow the item to be focused, with the remaining elements then appearing to the right of \ml{ആണ്}:

\begin{enumerate}[label=\textbf{\arabic*.}]
    \item \textbf{Normal Order (Subject in Focus):}\\
    \ml{എന്റെ വീട് കേരളത്തിലാണ്.}\\
    \textit{ente vīṭŭ kēraḷattil-āṇŭ.}\\
    ``My home is in Kerala.''
    
    \item \textbf{Cleft / Focused Order (Locative Phrase in Focus):}\\
    \ml{കേരളത്തിലാണ് എന്റെ വീട്.}\\
    \textit{kēraḷattil-āṇŭ ente vīṭŭ.}\\
    ``It's in Kerala that my home is.''\\
    \textbf{Focus:} \ml{കേരളത്തിൽ} (\textit{kēraḷattil} — ``in Kerala'').
\end{enumerate}

Because an equative sentence contains no lexical verb, no verbal nominalisation is triggered; focus is achieved purely through copular movement and constituent fronting.

\subsection{Focus on a Constituent in a Verbal Sentence}
As Moag \citep{moag2002malayalam} demonstrates, this same copular mechanism is extended to verbal sentences by, in effect, transforming them into equative cleft sentences. Here, \ml{ആണ്} is attached to the element to be brought into focus (usually, though not necessarily, appearing at the beginning of the sentence). The remaining elements, including the sentence verb, become the predicate for the newly created equative sentence. In order to fulfill the syntactic rule that the predicate contain a nominal, the original lexical verb of the sentence is converted into a noun by suffixing \ml{-അത്} (\textit{-atŭ}).

\subsubsection{Focus in Declarative Sentences}
In a declarative verbal sentence, any constituent (such as subject, direct object, or oblique argument) can be brought into focus:

\begin{enumerate}[label=\textbf{\arabic*.}]
    \item \textbf{Unfocused Baseline:}\\
    \ml{ഞാൻ അയാളെ അയച്ചു.}\\
    \textit{ñān ayāḷe ayaccu.}\\
    ``I sent him.''
    
    \item \textbf{Subject Focus:}\\
    \ml{ഞാനാണ് അയാളെ അയച്ചത്.}\\
    \textit{ñān-āṇŭ ayāḷe ayaccatŭ.}\\
    ``I am the one who sent him.'' (\textbf{Focus:} \ml{ഞാൻ} — ``I'').
    
    \item \textbf{Direct Object Focus:}\\
    \ml{ഞാൻ അയാളെയാണ് അയച്ചത്.}\\
    \textit{ñān ayāḷe-y-āṇŭ ayaccatŭ.}\\
    ``It's him that I sent.'' / ``He's the one that I sent.'' (\textbf{Focus:} \ml{അയാളെ} — ``him'').
    
    \item \textbf{Directional / Locative Focus:}\\
    \textbf{Neutral:} \ml{ഞാൻ കൊച്ചിക്ക് പോയി.} (\textit{ñān koccikkŭ pōyi} — ``I went to Cochin.'')\\
    \textbf{Cleft:} \ml{കൊച്ചിക്കാണ് ഞാൻ പോയത്.}\\
    \textit{koccikk-āṇŭ ñān pōyatŭ.}\\
    ``It is to Cochin that I went.'' (\textbf{Focus:} \ml{കൊച്ചിക്ക്} — ``to Cochin'').
\end{enumerate}

\subsubsection{Focus in Information Questions (Wh-Questions)}
In questions of the information type, only the question phrase can be brought into focus \citep{moag2002malayalam}. If the question consists of a single interrogative word such as \ml{ഏത്} (\textit{ētŭ} — ``which'') or \ml{ആര്} (\textit{ārŭ} — ``who''), \ml{ആണ്} attaches directly to the question word:

\begin{enumerate}[label=\textbf{\arabic*.}]
    \item \textbf{Unfocused Question:}\\
    \ml{നീ എന്ത് വാങ്ങിച്ചു?}\\
    \textit{nī entŭ vāṅṅiccu?}\\
    ``What did you buy?''
    
    \item \textbf{Focused Question:}\\
    \ml{നീ എന്താണ് വാങ്ങിച്ചത്?}\\
    \textit{nī ent-āṇŭ vāṅṅiccatŭ?}\\
    ``What is it that you bought?'' / ``What did you buy?''
\end{enumerate}

If the question is expressed by an entire phrase, \ml{ആണ്} attaches to the final word in that phrase:
\begin{quote}
\ml{ഏത് കടയിൽനിന്നാണ് വാങ്ങിച്ചത്?}\\
\textit{ētŭ kaṭayilninn-āṇŭ vāṅṅiccatŭ?}\\
``From which shop did you buy it?''
\end{quote}

\subsubsection{Negative and Contrastive Focus}
The same structural rule applies to phrases in declarative sentences whether positive or negative, providing a powerful vehicle for contrastive focus:
\begin{quote}
\ml{കേരളത്തിലല്ല കന്യാകുമാരി സ്ഥിതി ചെയ്യുന്നത്. തമിഴ് നാട്ടിലാണ് (സ്ഥിതി ചെയ്യുന്നത്).}\\
\textit{kēraḷattil-alla kanyākumāri sthiti ceyyunnatŭ. tamil nāṭṭil-āṇŭ (sthiti ceyyunnatŭ).}\\
``It's not in Kerala that Kanyakumari is situated. It's in Tamil Nadu (that it is situated).''
\end{quote}

\subsection{Focus on the Verb Itself}
\label{sec:focus_on_verb}
Moag \citep{moag2002malayalam} provides a dedicated treatment of focus on the verb itself under ``Cleft Sentences with Emphasis on the Verb'' (Sections 11.2 and 15.1). When the verbal action is the element being emphasized, two primary configurations are distinguished:

\subsubsection{Direct Finite Verb Emphasis}
The finite verb itself is nominalised and immediately followed by the copula \ml{ആണ്}:
\begin{quote}
\ml{ഞാൻ അയാളെ അയച്ചതാണ്.}\\
\textit{ñān ayāḷe ayaccat-āṇŭ.}\\
``I did send him.''
\end{quote}
Here, the truth-value or actual completion of the verbal action is emphasized, corresponding functionally to the emphatic auxiliary ``did'' in English.

\subsubsection{Emphatic Order on Purposive Action}
When an action is expressed by an infinitive or purposive verbal noun, emphatic cleft order places \ml{ആണ്} immediately after that verbal constituent in the middle of the sentence, while the backgrounded matrix verb appears as a nominalised verbal noun at the end:
\begin{quote}
\ml{ഞാൻ കൊച്ചിക്ക് പോകാനാണ് വന്നത്.}\\
\textit{ñān koccikkŭ pōkān-āṇŭ vannatŭ.}\\
``I (really) came in order to go to Cochin.''
\end{quote}
In this structure, the purposive action \ml{പോകാൻ} (\textit{pōkān} — ``to go'') is the focused element followed by \ml{ആണ്}, while \ml{വന്നത്} (\textit{vannatŭ} — ``the coming'') forms the presupposed predicate.

\subsubsection{Unemphatic / Softening Cleft Order}
Moag contrasts emphatic cleft order with what he terms the \textit{unemphatic order}, where the copula appears at the very end of the sentence, after the verbal noun:
\begin{enumerate}[label=\textbf{\arabic*.}]
    \item \ml{ഞാൻ എന്റെ മകനെയും ഭാര്യയെയും കൊച്ചിയിലേക്ക്}\\
    \ml{യാത്ര അയക്കാൻ വന്നതാണ്.}\\
    \textit{ñān ente makaneyuṁ bhāryayeyuṁ kocciyilēkkŭ yātra ayakkān vannat-āṇŭ.}\\
    ``I came to see my son and (his) wife off to Cochin.''
    
    \item \ml{ഞാൻ അവരോട് സംസാരിക്കാൻ പോയതാണ്.}\\
    \textit{ñān avarōṭŭ sansārikkān pōyat-āṇŭ.}\\
    ``I just went to talk with them.''
\end{enumerate}
Cleft sentences with this clause-final copular order are neutral with respect to emphasis and focus. Rather than highlighting a constituent, their use serves as a conversational softener, imparting a tone of offhandedness, politeness, or apology to the utterance.
